\section{Spectral modules, filter roots, and exact quantum cycles}
\label{sec:spectral-modules}
This section strengthens the temporal support theorem in five directions.
The classical statements concern observable covariance, not eigenvalue
statistics. The finite-field theorem is exact at every time, including times
of order the Hilbert-space dimension; it is not a local spectral law.

\subsection{All temporal translates form a free, stable module}
Retain $A,C,g,h_r$ and the isolated frequencies $k_r$ of
Lemma~\ref{lem:temporal-support}. On each oriented frequency orbit choose a
representative $k_{\mathcal O}$ and define the finite Laurent polynomial
\[
 H_r(\mathcal O,z)=\sum_{\ell\in\mathbb Z}
       \widehat h_r(A^\ell k_{\mathcal O})z^\ell .
\]
Only finitely many orbits and coefficients occur for fixed $m$. Let
$M_m(z)=(H_r(\mathcal O,z))_{\mathcal O,1\le r\le m}$ and
$S_m(\theta)=M_m(e^{i\theta})^*M_m(e^{i\theta})$.
Set
\[
 D_m^*=(2\pi^2)^m\prod_{r=1}^m
  \det(u_r,A^rv_r)^2|\widehat g(u_r)\widehat g(v_r)|^2,
 \qquad L_m=\sum_{r=1}^m\|\widehat h_r\|_1^2.
\]
Both are strictly positive.
\begin{theorem}[Uniform spectral positivity and temporal Riesz bounds]
\label{thm:module-riesz}
For all real $\theta$,
\[
 \det S_m(\theta)\ge D_m^*,\qquad
 \frac{D_m^*}{L_m^{m-1}}I_m\le S_m(\theta)\le L_m I_m.
\]
Consequently for every finitely supported complex coefficient array
$(b_{r,j})$, with arbitrary integer times $j$,
\begin{equation}\label{eq:module-riesz}
 \frac{D_m^*}{L_m^{m-1}}\sum_{r,j}|b_{r,j}|^2
 \le\left\|\sum_{r=1}^m\sum_j b_{r,j}C^jh_r\right\|_2^2
 \le L_m\sum_{r,j}|b_{r,j}|^2.
\end{equation}
The family $h_1,h_2,\ldots$ is free over $\mathbb C[z,z^{-1}]$ under
$z\cdot f=Cf$. In particular all functions
$\{C^ig,C^jg\}$, indexed by integer $i<j$, are linearly independent over
$\mathbb C$ in finite combinations.
\end{theorem}
\begin{proof}
Use $k_r$ and $-k_r$ themselves as the representatives of their pairwise
distinct orbits. In the rows of $M_m$ corresponding to this pair, columns
$j<r$ vanish, while column $r$ is the constant
$z_r=-\pi\det(u_r,A^rv_r)\widehat g(u_r)\widehat g(v_r)$ or its conjugate.
Selecting one row from each pair gives $2^m$ triangular minors with squared
absolute determinant $\prod_r|z_r|^2$, independently of $\theta$.
Cauchy--Binet proves the determinant bound. The triangle inequality gives
$\operatorname{Tr}S_m(\theta)\le L_m$, so every eigenvalue is at most $L_m$
and the smallest is at least $D_m^*/L_m^{m-1}$.

Transport by $C^j$ shifts the Fourier exponent on an orbit by $j$.
Parseval on each orbit identifies the middle member of
\eqref{eq:module-riesz} with
$\int\|M_m(e^{i\theta})b(e^{i\theta})\|^2d\theta/(2\pi)$,
where $b_r(z)=\sum_jb_{r,j}z^j$. A second application of Parseval proves the
bounds and module freeness. The last assertion uses
$\{C^ig,C^jg\}=4\pi C^ih_{j-i}$.
\end{proof}
The constants depend on the largest lag $m$ and the symbol $g$ but not on the
length or location of the time window. Thus the assertion does not contradict
the soft-mode theorem, which increases the number and size of the lags.
It also shows that no nonzero finite scalar combination of the $h_r$ is an
$L^2$ twisted coboundary $(C-e^{i\theta_0})F$: its orbit polynomial on the
isolated highest-lag orbit is a nonzero constant. Solving the corresponding
one-dimensional difference equation would require a nonsquare-summable tail.

\subsection{Complete coercivity classification for filtered single modes}
Let $g_0=a_0e_v+\overline{a_0}e_{-v}$ as before. Fix any nonzero real
polynomial $P(z)=\sum_{i=0}^dp_iz^i$, put $g=P(C)g_0$, and take $R\ge d$.
For $r=R+1,\ldots,R+n$, form the covariance matrix of
$\lambda^{-r}[h_r(g)]$, denoted $\widetilde G_{P,R,n}$.
Let $\nu_P(w)$ be the root multiplicity of $P$ at $w$ (zero if not a root),
and define
\[
 \mu_P(\theta)=\nu_P(\lambda e^{i\theta})+
                 \nu_P(\lambda^{-1}e^{-i\theta}),\qquad
 M_P=\max_{\theta\in\mathbb R}\mu_P(\theta).
\]
The maximum is zero when there are no roots on either indicated circle.
\begin{theorem}[Root-circle classification]\label{thm:root-circle}
With constants depending on $A,v,a_0,P,R$ but not on $n$,
\[
 \lambda_{\min}(\widetilde G_{P,R,n})\asymp n^{-2M_P}.
\]
In particular lag-uniform coercivity holds if and only if $P$ has no roots
on $|z|=\lambda$ or $|z|=\lambda^{-1}$.
\end{theorem}
\begin{proof}
Set $b_k=\sum_i p_ip_{i+k}$, extended evenly and by zero, and $w_k=b_k\lambda^k$.
The proof of \eqref{eq:soft-sandwich} applies unchanged and gives
\[
 C_0(1-\lambda^{-2(R+1-d)})^2T_n(F_P)
 \le\widetilde G_{P,R,n}\le C_0T_n(F_P),\qquad
 F_P(\theta)=|P(\lambda e^{i\theta})P(\lambda^{-1}e^{-i\theta})|^2.
\]
Here $C_0=4\pi^2|a_0|^4\det(v,Av)^2/(t^2-4)>0$.
The zero order of $F_P$ at $\theta_0$ is exactly $2\mu_P(\theta_0)$.
We give a direct finite-section estimate so that no external Toeplitz
asymptotic theorem is needed.

If there are no zeros, continuity and positivity bound $F_P$ above and below.
Otherwise let $Z$ be its finite zero set, of cardinality $J$, and $M=M_P$.
There is $c>0$ such that
$F_P(\theta)\ge c\,\operatorname{dist}(\theta,Z)^{2M}$, using circular
angular distance. For a degree-$(n-1)$ Fourier polynomial $p$ with normalized
$L^2$ norm one, $|p(\theta)|^2\le n$. The union of the radius
$\varepsilon=\pi/(4Jn)$ arcs about $Z$ therefore carries at most $1/4$ of
its squared mass. Outside these arcs $F_P\ge c\varepsilon^{2M}$.
This gives a lower bound $c'n^{-2M}$ for every Rayleigh quotient.

Choose a zero $\theta_0$ of maximal order and factor the finite Laurent
polynomial $P(\lambda z)P(\lambda^{-1}z^{-1})$ by
$(z-e^{i\theta_0})^M$. The remaining Laurent polynomial is bounded on the
unit circle. Use the coefficients
$x_j=e^{-ij\theta_0}\phi(j/(n+1))$, where
$\phi(x)=x^M(1-x)^M$ on $[0,1]$ and zero otherwise.
Its zero extension has bounded weak derivative of order $M$. The twisted
$M$th finite differences of $x$ have squared sum $O(n^{1-2M})$, while
$\|x\|^2\ge c_Mn$. Parseval gives the upper bound $O(n^{-2M})$.
For a real symmetric matrix a complex trial vector is harmless: its Rayleigh
quotient is a weighted average of those of its real and imaginary parts.
The quadratic-form comparison completes the proof.
\end{proof}
This specializes the familiar zero-order mechanism for Toeplitz eigenvalues
to an exact dynamical Gram representation. The general Toeplitz phenomenon
is prior theory, not a new universality mechanism; see B\"ottcher--Widom,
\emph{From Toeplitz Eigenvalues through Green's Kernels to Higher-Order
Wirtinger--Sobolev Inequalities}, arXiv:math/0412269.

There is also an exact repair for collective soft modes. Choose finitely many
filters $P_a$, take $R\ge\max_a\deg P_a$, and add their covariance matrices with the same
coefficient vector for their separate observations. Their symbols add.
Uniform coercivity is equivalent to absence of a common zero of the $F_{P_a}$.
If common zeros exist, the rate is $n^{-2M}$, where
\[
 M=\max_{\theta:\,F_{P_a}(\theta)=0\ \forall a}
                  \min_a\mu_{P_a}(\theta).
\]
The preceding lower and upper proofs apply to this positive sum; there can be
no cancellation of its leading nonnegative terms.
For example $P_\pm(z)=(z^2\mp tz+1)^q$ separately give $n^{-4q}$, but with
$u=2-2\cos\theta$ their sum obeys
\[
 F_+(\theta)+F_-(\theta)
 =u^{2q}(t^2-4+u)^{2q}
 +(4-u)^{2q}(t^2-u)^{2q}
 \ge 2^{2q+1}(t^2-4)^{2q}>0.
\]
Thus two explicit individually soft observations have a jointly uniform
positive lower covariance bound at every lag-window size. This statement
adds covariances of separate observations, not the covariance of their sum.

\subsection{The high-lag classification for every finite Fourier symbol}
The single-orbit restriction can be removed. Decompose the nonzero Fourier
support of an arbitrary real trigonometric polynomial into oriented $A$-orbits.
Choose one integer representative $w_a$ from each, shifting it so that
\[
 g=\sum_{a=1}^J P_a(C)e_{w_a},\qquad
 P_a(z)=\sum_{i=0}^{d_a}c_{a,i}z^i,\qquad P_a(0)\ne0.
\]
Representatives and filters can be paired by $w_{-a}=-w_a$ and
$P_{-a}=\overline{P_a}$ (coefficient conjugation). Put $d=\max_a d_a$ and
let $P_g$ be the monic greatest common divisor of the $P_a$ in $\mathbb C[z]$.
It has real coefficients, and changing representatives only multiplies filters
by Laurent monomials, which does not change its nonzero roots.

\begin{lemma}[Four-atom separation]\label{lem:four-atom-separation}
For a fixed finite nonzero integer alphabet $V=-V$, there is an integer $G$
such that every relation $\sum_{i=1}^4 A^{n_i}v_i=0$, $v_i\in V$, splits
into two vanishing subrelations across any consecutive time gap larger than $G$.
If $r,s>3G$ and
\[
 u+A^rv=A^j(u'+A^sv'),\qquad u,v,u',v'\in V,
\]
then necessarily $u=A^ju'$ and $A^rv=A^{j+s}v'$.
\end{lemma}
\begin{proof}
Set $B=\operatorname{sign}(\operatorname{tr}A)A$, and choose its eigen-covectors
$\xi,\eta$. There is $\kappa\ne0$ with
$Q_B(x)=\det(x,Bx)=\kappa\xi(x)\eta(x)$.
Let $X=\max_V|\xi|$, $Y=\max_V|\eta|$, and choose $G$ so that
$16|\kappa|XY\lambda^{-G}<1$.
Suppose a gap separates times at most $a$ from times at least $b$, $b-a>G$.
The sum $w$ of the terms on the left is integral, and, using minus the sum
on the right for its contracting coordinate,
\[
 |\xi(w)|\le4X\lambda^a,\qquad |\eta(w)|\le4Y\lambda^{-b}.
\]
Thus the integer $Q_B(w)$ has absolute value less than one. It vanishes,
and irrationality of the eigenlines implies $w=0$. The other subrelation
also vanishes. Signs from negative eigenvalues are harmless in these bounds.
The same argument shows that cancellation of two alphabet atoms can occur
only when their time difference has absolute value at most $G$.

For the second assertion, order the four times $0,r,j,j+s$. Their span
exceeds $3G$, so one of the three gaps exceeds $G$. Neither side can be a
singleton because no atom is zero. Hence they split into two cancelling pairs.
A pair cannot consist of the two atoms of the same side of the original
identity, since $r,s>3G$. Crossed pairing would require
$|j+s|\le G$ and $|r-j|\le G$, hence $r+s\le2G$, also impossible.
Only the asserted parallel pairing remains.
\end{proof}

For definiteness the constant in this lemma can be bounded using integers.
Write $B=(b_{ij})$, $t=\operatorname{tr}B$, $\Delta=t^2-4$.
The covectors
$\xi(x,y)=2b_{21}x+(t-2b_{11}+\sqrt\Delta)y$ and its conjugate satisfy
$\xi\eta=4b_{21}Q_B$; here $b_{21}\ne0$.
Taking
$M=\max_{(x,y)\in V}\{2|b_{21}||x|+(|t-2b_{11}|+t)|y|\}$,
any integer $G$ with $2^G>4M^2/|b_{21}|$ suffices, since $\lambda>2$.

Choose eigenvectors $e_+,e_-$ of $B$ and write
$w_a=\xi_a e_++\eta_a e_-$. Every $\xi_a,\eta_a$ is nonzero.
Set $D_e=\det(e_+,e_-)$ and
\[
 F_g(\theta)=\pi^2 D_e^2
 \left(\sum_a\eta_a^2|P_a(\lambda^{-1}e^{-i\theta})|^2\right)
 \left(\sum_a\xi_a^2|P_a(\lambda e^{i\theta})|^2\right).
\]
This expression is independent of eigenvector rescaling.
\begin{theorem}[Complete finite-symbol high-lag criterion]
\label{thm:general-symbol}
There is $R_0(A,g)$ such that, for every fixed $R\ge R_0$, the covariance
matrix $\widetilde G_{g,R,n}$ of $\lambda^{-r}[h_r(g)]$,
$r=R+1,\ldots,R+n$, satisfies
\[
 c\,T_n(F_g)\le\widetilde G_{g,R,n}\le C\,T_n(F_g)
 \quad(n\ge1)
\]
with positive constants independent of $n$. Therefore
\[
 \lambda_{\min}(\widetilde G_{g,R,n})\asymp n^{-2M_g},\qquad
 M_g=\max_\theta\{
 \nu_{P_g}(\lambda e^{i\theta})+
 \nu_{P_g}(\lambda^{-1}e^{-i\theta})\}.
\]
In particular uniform high-lag coercivity holds precisely when the common
finite temporal divisor $P_g$ has no root on the two hyperbolic circles.
This covers every finite real Fourier symbol, with arbitrary phases and any
number of frequency orbits.
\end{theorem}
\begin{proof}
Apply Lemma~\ref{lem:four-atom-separation} to the actual Fourier support $V$.
For sufficiently large positive $r,s$, equality of two contributing frequency
orbits is possible only by parallel matching of their two atoms. In the
chosen orbit coordinates, this says that the orbit is uniquely indexed by
an ordered pair $(a,b)$ and an effective lag $\ell=r+j-i$. Thus its Fourier
sum is
\[
 -\pi D_{ab}(\ell)\,b_{ab}(\ell-r),\qquad
 D_{ab}(\ell)=\det(w_a,A^\ell w_b),\quad
 b_{ab}(k)=\sum_i c_{a,i}c_{b,i+k}.
\]
The coefficients have been extended by zero; $|k|\le d$.
Consequently the exact covariance, before normalization, is
\[
 (G_{g,R,n})_{rs}
 =\pi^2\sum_{a,b,\ell}D_{ab}(\ell)^2
        b_{ab}(\ell-r)\overline{b_{ab}(\ell-s)}.
\]
If $s_A=\operatorname{sign}(\operatorname{tr}A)$, diagonalization gives
\[
 D_{ab}(\ell)=s_A^\ell
 \{-D_e\eta_a\xi_b\lambda^\ell+D_e\xi_a\eta_b\lambda^{-\ell}\}.
\]
The leading coefficient $\gamma_{ab}=-D_e\eta_a\xi_b$ is nonzero.
Choose $R_0$ so that all $r,s>3G$, all $\ell\ge R+1-d>0$, and
\[
 \left|\frac{D_e\xi_a\eta_b}{\gamma_{ab}}\right|
          \lambda^{-2\ell}\le\tfrac12
 \quad\hbox{for every }a,b.
\]
Then the squared residual factors are between $1/4$ and $9/4$.
For $w_{ab,k}=b_{ab}(k)\lambda^k$, the normalized quadratic form is a sum
of their convolution norms with these positive residual weights. Parseval
gives the claimed comparison with constants $1/4$ and $9/4$ and symbol
\[
 \pi^2\sum_{a,b}\gamma_{ab}^2
  |P_a(\lambda^{-1}e^{-i\theta})P_b(\lambda e^{i\theta})|^2=F_g(\theta).
\]
There is no missing conjugation in $b_{ab}$: it arises from the bilinear
Poisson bracket. Conjugation enters only in the covariance and absolute square.
The reality pairs make $F_g$ even and the covariance real.

Each positive sum in $F_g$ vanishes precisely when all of its filters vanish.
Its vanishing multiplicity is twice their minimum root multiplicity, namely
that of their gcd. The product therefore has zero order $2M_g$ at its
most degenerate zero. The multi-zero finite-section estimate in
Theorem~\ref{thm:root-circle} applies to this symbol as well: it is a finite
nonnegative trigonometric polynomial; near a maximal zero it is bounded above
by a constant times $|e^{i\theta}-e^{i\theta_0}|^{2M_g}$.
The same concentration lower bound and modulated finite-difference trial
vector prove the asserted rate.
\end{proof}
For example, if one of the orbit filters is a nonzero constant, then
$P_g=1$ and the high-lag covariance is uniformly coercive. More generally,
adding a second orbit with a coprime temporal filter removes the common-root
obstruction within a single observable. This is stronger than merely adding
covariance matrices of separate observations, but relies on the exact
large-lag orbit separation just proved. No claim about infinite Fourier
support or nonlinear kicked spectral statistics follows from this criterion.


\begin{corollary}[One-harmonic repair and its crossover scale]
\label{cor:one-harmonic-repair}
Let $v\ne0$ be integral, $g_v=e_v+e_{-v}$, $q\ge1$, and
\[
 g_\varepsilon=(C^2-tC+I)^qg_v+\varepsilon(e_{2v}+e_{-2v}),
 \qquad 0\le\varepsilon\le1.
\]
For a sufficiently large fixed starting lag $R$, there are positive constants
$c,C$, independent of both $n$ and $\varepsilon$, such that
\[
 c\,(n^{-2q}+\varepsilon^2)^2
 \le\lambda_{\min}(\widetilde G_{g_\varepsilon,R,n})
 \le C\,(n^{-2q}+\varepsilon^2)^2.
\]
Thus a single additional harmonic changes the arbitrarily soft $n^{-4q}$
family into a uniformly coercive one. The crossover occurs at
$n\asymp\varepsilon^{-1/q}$, and the large-window floor has order
$\varepsilon^4$.
\end{corollary}
\begin{proof}
The oriented orbits of $v$ and $2v$ are distinct, as
$Q_B(2v)=4Q_B(v)\ne Q_B(v)$. The four filters in
Theorem~\ref{thm:general-symbol} are $P_q,P_q,\varepsilon,\varepsilon$;
the same separation and residual bounds work for all $\varepsilon$, including
zero, by using the larger fixed support. Set $u=2-2\cos\theta$,
$B(\theta)=u(t^2-4+u)$ and
$C_0=4\pi^2\det(v,Av)^2/(t^2-4)$. The exact Toeplitz symbol is
\[
 F_{g_\varepsilon}(\theta)
 =C_0\{B(\theta)^{2q}
       +4\varepsilon^2(\lambda^{2q}+\lambda^{-2q})B(\theta)^q
       +16\varepsilon^4\}.
\]
Since $t^2-4>0$ and $0\le u\le4$, this is bounded above and below by
positive constant multiples of $(u^q+\varepsilon^2)^2$, uniformly in
$\varepsilon$. For every normalized degree-$(n-1)$ polynomial $p$, Jensen's
inequality bounds its integrals against $u^q$ and $u^{2q}$ below by the
corresponding powers of the Dirichlet Laplacian minimum
$4\sin^2(\pi/(2(n+1)))$. This gives the required lower estimate.
The trial coefficients from the soft-mode theorem have normalized $q$th and
$2q$th finite-difference squared norms $O(n^{-2q})$ and $O(n^{-4q})$.
They yield the upper estimate including both cross and constant terms.
The exact finite-covariance comparison completes the proof.
\end{proof}
This repair is a theorem about the covariance of one perturbed classical
observable. It does not claim that this harmonic perturbation proves BGS for
the corresponding quantum propagator.

\subsection{An exact all-time quantum module over a finite field}
Let $p$ be an odd prime and $A\in\mathrm{SL}(2,\mathbb F_p)$ satisfy
$(\operatorname{tr}A)^2\ne4$. Assume its exact order $K\ge3$ is odd.
Choose $v\in\mathbb F_p^2$ with $\det(v,Av)\ne0$.
Use modular-centred Weyl matrices $W_p(w)$, so
\begin{align*}
 W_p(u)W_p(w)&=e_p(\omega(u,w)/2)W_p(u+w),&
 W_p(w)^*&=W_p(-w),\\
 p^{-1}\operatorname{Tr}(W_p(u)^*W_p(w))&=\mathbf1_{u=w}.
\end{align*}
Here $1/2$ is the inverse of two in $\mathbb F_p$, not ordinary division.
The map $\alpha(W_p(w))=W_p(Aw)$ extends to a $*$-automorphism of the full
matrix algebra, hence is implemented by a unitary. This fixes the finite
quantization convention; it does not assert physical-lift Egorov in every
boundary sector for every integral matrix.
Put $G=cW_p(v)+\overline cW_p(-v)$, $c\ne0$, and
\[
 H_r=\frac{ip}{2}[G,\alpha^rG],\quad
 1\le r\le m=(K-1)/2,\quad d_r=\det(v,A^rv)\in\mathbb F_p.
\]
Write $\widetilde d_r$ for any integer representative of $d_r$ and
$\tau_p=p^{-1}\operatorname{Tr}$.
\begin{theorem}[Cyclic module and variance at every time]
\label{thm:quantum-cycles}
For any finitely supported complex coefficients $b_{r,j}$, $j\in\mathbb Z$,
\begin{equation}\label{eq:quantum-cyclic-module}
 \tau_p\left|\sum_{r=1}^m\sum_j b_{r,j}\alpha^jH_r\right|^2
 =4p^2|c|^4\sum_{r=1}^m\sin^2(\pi\widetilde d_r/p)
       \sum_{a=0}^{K-1}\left|\sum_{j\equiv a\bmod K}b_{r,j}\right|^2.
\end{equation}
All $d_r$ are nonzero. The $Km$ operators
$\alpha^aH_r$, $0\le a<K$, $1\le r\le m$, are pairwise orthogonal and
nonzero. The kernel of the temporal Laurent-polynomial map is exactly
$(z^K-1)\mathbb C[z,z^{-1}]^m$.

For $n=qK+s$, $0\le s<K$, in particular,
\begin{equation}\label{eq:quantum-all-time}
 \tau_p\left|\sum_{j=0}^{n-1}\alpha^jH_r\right|^2
 =4p^2|c|^4\sin^2(\pi\widetilde d_r/p)
       \left\{\frac{n^2}{K}+s\left(1-\frac{s}{K}\right)\right\}.
\end{equation}
This holds for every integer $n\ge1$, without a no-alias time restriction.
\end{theorem}
\begin{proof}
Over a splitting field, $A$ has distinct eigenvalues $\zeta,\zeta^{-1}$ of
exact odd order $K$. Every nonzero vector has orbit length $K$, and no such
orbit contains both $w$ and $-w$: otherwise a power of $\zeta$ would be $-1$.
Let $Q(w)=\det(w,Aw)$ and $t_r=\zeta^r+\zeta^{-r}$. Then
\[
 Q(v\pm A^rv)=(2\pm t_r)Q(v).
\]
For $1\le r\le(K-1)/2$, all the $2m$ values $2\pm t_r$ are nonzero and
distinct. Equality $t_r=t_s$ implies $r\equiv\pm s\pmod K$, while
$t_r=-t_s$ would imply $\zeta^r=-\zeta^{\pm s}$, impossible at odd order.
Thus the $4m$ frequencies $\pm(v+A^rv)$, $\pm(v-A^rv)$ lie in pairwise
distinct orbits. The condition $Q(v)\ne0$ also gives
$d_r=Q(v)(\zeta^r-\zeta^{-r})/(\zeta-\zeta^{-1})\ne0$.

Expanding the Weyl commutator shows that $H_r$ has exactly these four
frequencies, each with coefficient modulus
$p|c|^2|\sin(\pi\widetilde d_r/p)|$. The modular phase may differ by a
sign from ordinary half-division, but its squared modulus is the displayed
one. Transport permutes each orbit without an additional phase in this
convention. Since the orbits are disjoint, collecting coefficients by their
residue class in time and applying Weyl orthogonality proves
\eqref{eq:quantum-cyclic-module}. The only polynomial kernel is reduction to
zero modulo $z^K-1$. For the constant time coefficients there are $s$ residue
classes visited $q+1$ times and $K-s$ classes visited $q$ times; their squared
counts sum to $Kq^2+(2q+1)s=n^2/K+s(1-s/K)$.
\end{proof}
For $n\le K$ the brace equals $n$, giving exact linear variance. For $n\gg K$
it is $n^2/K+O(K)$, exposing coherent recurrences. When $n$ is of order $p$
the formula still holds, regardless of whether $K$ is comparable to $p$;
linear variance at that scale additionally requires $K\ge n$.
No infinite sequence with $K\asymp p$ for one fixed integral $A$ is inferred.
The theorem is about an unperturbed finite quantum map and commutator
observables, not the local eigenphase statistics of a kicked propagator.
The odd-order hypothesis is essential: at $p=7$ and
$A=\left(\begin{smallmatrix}3&-1\\1&0\end{smallmatrix}\right)$ the order is
$8$ and opposite frequencies share an orbit, invalidating the displayed
orthogonality and variance formula. This is an explicit negative control.


\subsection{Phase completion at even order and an exact rank transition}
The failure at even order is not unavoidable. It depends on the Fourier phase
of the real initial observable, and can be removed without randomness.
Retain the regular-semisimple assumptions over $\mathbb F_p$, but allow any
order $K\ge3$, and put $m=\lfloor(K-1)/2\rfloor$.
\begin{theorem}[Balanced-phase completion]\label{thm:phase-completion}
If $K$ is even, assume $c=a+ib$ with $a,b\in\mathbb R$ and $a^2=b^2>0$.
Then \eqref{eq:quantum-cyclic-module} and \eqref{eq:quantum-all-time} hold
unchanged for all $1\le r\le m$. Thus the cyclic module is free of rank $m$
over $\mathbb C[z,z^{-1}]/(z^K-1)$ for every regular semisimple order,
using an arbitrary nonzero phase at odd order and a balanced phase at even order.

More generally, at even order let $ab\ne0$, set
$s_r=\sin(2\pi \overline2\widetilde d_r/p)$, where
$\overline2=(p+1)/2$, and define $J_r=H_r/(ps_r)$. If
$B_{r,u}=\sum_{j\equiv u\bmod K}b_{r,j}$, then
\[
 16\min(a^4,b^4)\sum_{r,u}|B_{r,u}|^2
 \le\tau_p\left|\sum_{r,j}b_{r,j}\alpha^jJ_r\right|^2
 \le16\max(a^4,b^4)\sum_{r,u}|B_{r,u}|^2.
\]
The dimension of the span of all $\alpha^jH_r$ is $K(K-2)/2$ for $ab\ne0$,
but is exactly $K(K-2)/8$ for a nonzero purely real or purely imaginary $c$.
\end{theorem}
\begin{proof}
Write $K=2L$. Then $A^L=-I$. The invariant $Q(w)=\det(w,Aw)$ groups
the lags $r$ and $L-r$, for $1\le r<L/2$, into pairs. Different such
pairs have disjoint label orbits. Indeed $t_r=t_s$ implies $r=\pm s$
modulo $K$, while $t_r=-t_s$ now implies $r=L\pm s$ modulo $K$.
Within a pair choose $E=W_p(v+A^rv)$ and $F=W_p(v-A^rv)$ as orbit seeds.
Their $K$-element orbits are disjoint. Expanding the commutators gives, in
the cyclic group algebra, with $z$ acting by $\alpha$,
\[
 J_r=-A_c(z)E+B_c(z)F,\qquad
 z^rJ_{L-r}=B_c(z)E-A'_c(z)F,
\]
where
\[
 A_c(z)=c^2+\overline c^{\,2}z^L,\qquad
 A'_c(z)=\overline c^{\,2}+c^2z^L,\qquad
 B_c(z)=|c|^2(1+z^L).
\]
Here $d_{L-r}=d_r$, so the two sine normalizations agree. The frequency
identities used are
$v+A^{L-r}v=-A^{-r}(v-A^rv)$ and
$v-A^{L-r}v=A^{-r}(v+A^rv)$.

A unitary Fourier transform on each cyclic orbit reduces the block to
\[
 \begin{pmatrix}-A_c(z)&B_c(z)\\ B_c(z)&-A'_c(z)\end{pmatrix},
 \qquad z^K=1.
\]
For $z^L=1$, its singular values are $4a^2$ and $4b^2$.
For $z^L=-1$, both singular values are $4|ab|$.
Multiplication of the second generator by $z^r$ is unitary and does not
change a bound on coefficient norms.

If $L$ is even, the remaining lag $r=L/2$ is a separate block: all four
frequencies are on the single orbit of $E=W_p(v+A^rv)$, since
$v-A^rv=A^{-r}(v+A^rv)$. Its orbit polynomial is
\[
 -c^2-\overline c^{\,2}z^{2r}+|c|^2(z^r+z^{-r}).
\]
At $z^r=1,-1,i,-i$ its values are respectively
$4b^2,-4a^2,-4iab,-4iab$. Its squared singular values therefore satisfy
the same bounds. These blocks exhaust the half-period lags.
When $a^2=b^2$, every block is a scalar multiple of a unitary, with squared
singular value $16a^4=4|c|^4$. Substituting $ps_rb_{r,j}$ for coefficients
of $J_r$ proves the claimed exact formula for $H_r$.

For $ab\ne0$ all blocks are invertible at all characters, giving $Km=2L(L-1)$.
For a purely real or imaginary phase each two-generator block has rank one
on the $L$ characters with $z^L=1$ and rank zero on the others. Its total
rank is $L$. If present, the exceptional one-generator block has rank $K/4$.
Counting $(L-1)/2$ pairs when $L$ is odd, or $(L-2)/2$ pairs and one
exceptional block when $L$ is even, gives $L(L-1)/2=K(K-2)/8$.
\end{proof}
This is a phase-controlled rank transition in the module of commutator
observables. It is not a transition theorem for eigenvalue statistics.
The purely real even-order negative control remains valid; the balanced
choice $c=1+i$ removes precisely that phase-dependent degeneracy.
